\documentclass[reprint,amsmath,amssymb,aps,twoside]{revtex4-2} % for editor use only

% For initial submission use these lines
%\documentclass[amsmath,amssymb,aps,endfloats,linenumbers]{revtex4-2}



% DO NOT CHANGE THINGS BELOW HERE
\usepackage{graphicx}
\usepackage{amsmath,amssymb,amsfonts}
\usepackage{dcolumn}
\usepackage{bm}
\usepackage{siunitx}
%\usepackage{tikz,pgfplots}
\sisetup{separate-uncertainty=true, multi-part-units=single}
\usepackage[colorlinks,allcolors=blue]{hyperref}
\usepackage{cleveref}
\crefname{equation}{}{}
\crefname{figure}{Fig.}{Figs.}
\crefname{table}{Table}{Tables}
\usepackage{svg}
\svgpath{{./figures}}




% set PDF metadata
% AUTHORS EDIT THIS PART
\hypersetup{%
pdftitle={The effect of popper design on maximum jump height}, % full title here
pdfauthor={Kaitlin Coulanges, Saanvi Dakwale, Samantha Hein, Reese Wallace, and Sashank Yellapragada}, % full author list here
}


\usepackage{fancyhdr}
\pagestyle{fancy}
\fancyhf{}
\fancyhead[RE,RO]{J S\&E \textbf{2}, xx--xx (2026)} % FOR EDITOR ONLY
\fancyhead[LO]{Yellapragada \emph{et al}} % if many authors
%\fancyhead[LO]{John and Dee} % if short enough to list all authors
\fancyhead[LE]{The effect of popper design on maximum jump height} % full or short title depending on what fits
%\fancyhead[LE}{Several species of small furry animals gathered together in a cave and grooving with a pict} % might be too long
%\fancyhead[LE]{Several species grooving} % might be a short version of the title that fits
\fancyfoot[C]{\thepage}
\fancypagestyle{mytitlepage}{
\fancyhf{}
\fancyhead[C]{Journal of Science \& Engineering \textbf{2}, xx--xx (2026)} % FOR EDITOR ONLY
\fancyfoot[C]{\thepage}
%\fancyfoot[L]{\scriptsize\copyright\ 2026 Journal of Science \& Engineering, \href{https://creativecommons.org/licenses/by-nc/4.0/}{CC-BY-NC 4.0}}
}


\begin{document}
%\setcounter{page}{67} % FOR EDITOR USE ONLY
%\linenumbers % FOR EDITOR USE ONLY

% AUTHORS EDIT THIS PART
\title{The effect of popper design on maximum jump height}
\author{Kaitlin Coulanges}
\author{Saanvi Dakwale}
\author{Samantha Hein}
\author{Reese Wallace}
\author{Sashank Yellapragada}
\email[Contact author: ]{427syellapragada@frhsd.com} % use FRHSD address
% For S&E authors
\affiliation{Science \& Engineering Magnet Program, \href{https://manalapan.frhsd.com/}{Manalapan High School}, Englishtown, NJ 07726 USA}
\date{\today}


% AUTHORS PUT YOUR MANUSCRIPT IN BELOW HERE
\begin{abstract}
This experiment investigates whether the facial expression of a popper affects its maximum jump height using principles of energy conservation. Five ``happy'' poppers and five ``sad'' poppers were compressed to maximum depth and released, converting elastic potential energy into kinetic energy, and then gravitational potential energy. Elastic potential energy was measured via a force-compression graph ($EPE = \qty{294.0}{\milli\joule}$). Mean kinetic energy ($KE$) at launch was \qty{94.1}{\milli\joule} (happy) and \qty{81.4}{\milli\joule} (sad), representing approximately \num{32}\% and \num{28}\% of stored $EPE$ respectively. Paired $t$-tests confirmed that $KE$ at launch and $GPE$ at peak were not significantly different within either group (happy: $p = 0.41$; sad: $p = 0.43$), verifying that mechanical energy is conserved during flight. Two-sample $t$-tests found no statistically significant difference in mass ($p = 0.12$) or height ($p = 0.24$). These results support the hypothesis that facial expression does not affect energy transfer, and that popper motion is governed by elastic properties and mass rather than surface design.

%\vspace{1em} % Add space before DOI
%\noindent DOI: \href{https://doi.org/10.64804/xxxxxx}{10.64808/xxxxxx}
\end{abstract}

\keywords{kinematics, energy conservation, elastic potential energy, poppers}

\maketitle\thispagestyle{mytitlepage}

%The current column width is \the\columnwidth






\section{Introduction} 
When a popper is compressed, work is done to deform the material, storing elastic potential energy ($EPE$). This work can be described by the area under the force versus compression graph. 
Upon release of the popper, $EPE$ is rapidly converted into kinetic energy ($KE$), given by:
\begin{equation} 
KE = \frac{1}{2} m v^2,
\label{eq:ke}
\end{equation}
where $m$ is the mass and $v$ is the velocity \cite{tipler}. As the popper rises, $KE$ is transformed into gravitational potential energy (GPE). 
At maximum height ($h$), where velocity is approximately zero, GPE is given by:
\begin{equation}
GPE = mgh,
\label{eq:gpe}
\end{equation}
where $g=\qty{9.8}{\meter\per\second\squared}$ is the acceleration due to gravity \cite{tipler}.
If energy is conserved between launch and peak, then $KE$ will be approximately equal to $GPE$. In this experiment, we aim to determine whether the facial expression of a rubber popper affects its maximum jump height. Because the two designs share the same elastic properties and composition, and primarily only differ in their facial expressions, we hypothesize that facial expression will not significantly affect maximum height achieved. Additionally, we wish to find whether or not energy is conserved during the flight of the popper, acting as justification for using height as an indicator of energy transfer.







\section{Methods and materials} 
This experiment was conducted using commercially available rubber poppers (Liberty Imports; Shantou, CN) with multiple facial expressions. Two types were chosen: ``happy,'' identified by an upward-curved mouth, and ``sad,'' denoted by a downward-curved mouth and tear-like markings beneath each eye. A sample size of $n=5$ was used for each facial expression group. The mass of each popper was measured on a digital scale (Smart Weight; Zhongshan, China). 

All trials were conducted indoors on a flat, level surface within a classroom to maintain consistent environmental conditions. For measuring jump height, two meter sticks (Westcott Rule Company; Shelton, CT) were taped vertically to the wall. For each trial, to reduce experimental error, the popper was compressed to its maximum depth on the floor. Additionally, one student launched all the poppers. To ensure accuracy, the maximum height of each jump was documented on an iPhone 16 (Apple Inc; Cupertino, CA), recording at \qty{240}{fps} with \num{1080}p and a standard wide-angle lens, placed at a fixed height above the ground. 

Specific landmarks were used to calibrate height of the popper after launch, from which jump heights for each popper were determined and recorded. Launch velocity was measured using FizziQ \cite{fizziq}. The frame closest to the release of the popper was defined as $y = 0, t = 0$. To reduce bias, trial order was randomized rather than testing all poppers of one facial expression and trials in which the popper failed to launch properly were discarded. Outliers, defined as trials where the height was two or more standard deviations from the mean, were also excluded from the final analysis. To measure $EPE$, a digital scale was set to zero beneath the popper and a scissor jack was placed directly above it, compressing the popper in controlled increments. The mass readings were converted to force ($F = mg$) at each measured compression distance. $EPE$ was calculated as the area under the force-compression curve. $KE$ and $GPE$ were calculated from \cref{eq:ke,eq:gpe} for each trial. 
%The following hypotheses were tested using two-sampled, two-tailed $t$-tests ($\alpha=0.05$):
%Regarding masses:
%𝐻0, 𝐴: µℎ𝑎𝑝𝑝𝑦 𝑚𝑎𝑠𝑠 = µ𝑠𝑎𝑑 𝑚𝑎𝑠𝑠 , or, no difference in the mean masses of the poppers
%𝐻1, 𝐴: µℎ𝑎𝑝𝑝𝑦 𝑚𝑎𝑠𝑠 ≠ µ𝑠𝑎𝑑 𝑚𝑎𝑠𝑠 , or, a significant difference exists between the mean masses of the
%poppers
%Regarding heights:
%𝐻0, 𝐵: µℎ𝑎𝑝𝑝𝑦 ℎ𝑒𝑖𝑔ℎ𝑡 = µ𝑠𝑎𝑑 ℎ𝑒𝑖𝑔ℎ𝑡 , or, no difference in the mean jump height
%𝐻1, 𝐵: µℎ𝑎𝑝𝑝𝑦 ℎ𝑒𝑖𝑔ℎ𝑡 ≠ µ𝑠𝑎𝑑 ℎ𝑒𝑖𝑔ℎ𝑡 , or, a significant difference exists between the mean jump
%heights
% Additionally, paired two-tailed $t$-tests were performed to evaluate energy conservation:
%𝐻0, C: µdiff = 0 , (where $KE = GPE$)
%𝐻1, C: µdiff ≠ 0

Data analysis and the force-compression graph were done in Google Sheets \cite{googlesheets}. \texttt{R} with library \texttt{ggplot2} \cite{R, ggplot2}, was used to create a bar chart comparing the mean maximum jump height for happy and sad poppers.

The experimental setup is illustrated in \cref{fig:setup}. The setup for measuring $EPE$ is illustrated in \cref{fig:setup_epe}.
\begin{figure}
\begin{center}
\includesvg[width=\columnwidth]{figures/setupphoto.svg}
\end{center}
\caption{Experimental setup showing the vertical calibration scale and popper placement.}
\label{fig:setup} 
\end{figure}

\begin{figure}
\begin{center}
\includesvg[width=\columnwidth]{figures/scissor_jack.svg}
\end{center}
\caption{The scissor jack was used to apply a controlled force and the digital scale recorded the normal force at various compressions.}
\label{fig:setup_epe}
\end{figure}





\section{Results} 
The measured mass, maximum jump height, and calculated gravitational potential energy (GPE) for the happy and sad popper groups are summarized in \cref{tab:happy} and \cref{tab:sad}, respectively.
\begin{figure}
\begin{center}
\includesvg[width=\columnwidth]{figures/force_graph.svg}
\end{center}
\caption{Force versus compression graph used to calculate total $EPE$ via area under the curve.}
\label{fig:force_graph} 
\end{figure}
The force-compression graph is seen in \cref{fig:force_graph}. Total ($EPE$) was calculated by integrating the area under the curve in three segments, giving \qty{294.0}{\milli\joule}.

\begin{table}
\caption{Mass $m$, height $h$, and energy data for Happy Poppers.}
\label{tab:happy} 
\begin{ruledtabular}
\begin{tabular}{cS[table-format=1.4]S[table-format=1.2]S[table-format=3.1]} 
trial & {$m$, \unit{\kilo\gram}} & {$h$, \unit{\meter}} & {$v_{\text{launch}}$, \unit{\meter\per\second}} & {$KE$, \unit{\milli\joule}} & {$GPE$, \unit{\milli\joule}} & {KE - GPE, \unit{\milli\joule}} \\
\colrule \\
1 & 0.0058 & 2.12 & 6.26 & 113.6 & 120.6 & 7.0 \\
2 & 0.0057 & 1.82 & 6.10 & 105.9 & 101.8 & 4.1 \\
3 & 0.0056 & 1.47 & 5.16 & 74.4 & 80.8 & 6.3 \\
4 & 0.0057 & 1.67 & 5.78 & 95.3 & 93.4 & 1.9 \\
5 & 0.0056 & 1.54 & 5.39 & 81.3 & 84.6 & 3.4 \\
\colrule 
mean & 0.0057 & 1.72 & 5.74 & 94.1 & 96.2 & 4.6 \\
sd & 0.00008 & 0.26 & 0.46 & 16.3 & 15.9 & 2.2 \\
\end{tabular} 
\end{ruledtabular}
\end{table}

\begin{table}
\caption{Mass $m$, height $h$, and energy data for sad poppers.}
\label{tab:sad} 
\begin{ruledtabular}
\begin{tabular}{cS[table-format=1.4]S[table-format=1.2]S[table-format=2.1]} 
trial & {$m$, \unit{\kilo\gram}} & {$h$, \unit{\meter}} & {$v_{\text{launch}}$, \unit{\meter\per\second}} & {$KE$, \unit{\milli\joule}} & {$GPE$, \unit{\milli\joule}} & {KE - GPE, \unit{\milli\joule}} \\
\colrule \\
1 & 0.0050 & 1.63 & 5.77 & 83.2 & 80.0 & 3.2 \\
2 & 0.0055 & 1.30 & 4.90 & 66.0 & 70.1 & 4.1 \\
3 & 0.0056 & 1.53 & 5.42 & 82.4 & 84.1 & 1.7 \\
4 & 0.0057 & 1.72 & 5.58 & 88.6 & 96.2 & 7.5 \\
5 & 0.0055 & 1.58 & 5.62 & 87.0 & 85.2 & 1.7 \\
\colrule 
mean & 0.0055 & 1.55 & 5.46 & 81.4 & 83.1 & 3.7 \\
sd & 0.0002 & 0.16 & 0.34 & 9.0 & 9.4 & 2.4 \\
\end{tabular} 
\end{ruledtabular}
\end{table}

The results of the two-sample $t$-tests are presented in \cref{tab:stats}. At the $\alpha = 0.05$ significance level, no statistically significant differences were found between the two groups for either of the measured variables.
\begin{table}
\caption{Two-sample $t$-test results for mass $m$ and height $h$}
\label{tab:stats} 
\begin{ruledtabular}
\begin{tabular}{lSSSS} 
variable   & {$t$-value} & {df} & {$p$-value} \\
\colrule
mass $m$   & 1.74 & 8 & 0.12 \\
height $h$ & 1.27 & 8 & 0.24 \\
\end{tabular}
\end{ruledtabular} 
\end{table}

\begin{table} 
\caption{Paired $t$-test results for $KE$ at launch and $GPE$ at peak}
\label{tab:ttest_paired}
\begin{ruledtabular}
\begin{tabular}{lccc}
group & $t$-statistic & $df$ & $p$-value \\ 
\colrule
Happy & -0.93 & 4 & 0.41 \\ 
Sad & -0.88 & 4 & 0.43 \\
\end{tabular}
\end{ruledtabular}
\end{table}

\begin{figure}
\begin{center}
\includesvg[width=\columnwidth]{fig2.svg}
\end{center} 
\caption{Mean maximum jump height for happy and sad poppers. Error bars represent $\pm 1$ standard deviation ($n = 5$ per group).}
\label{fig:2} 
\end{figure}













\section{Discussion}
The results of this experiment indicate that popper design does not have a statistically significant effect on maximum jump height. Although the mean height of the happy poppers (\qty{1.72}{\meter}) was slightly greater than that of the sad poppers (\qty{1.55}{\meter}), the distributions overlap substantially, as indicated by \cref{fig:2}. \cref{tab:ttest_two} shows that $p=0.24$ for height, which is greater than $\alpha=0.05$. Therefore, we fail to reject the null hypothesis. Additionally, no significant difference was found in mass ($p = 0.12$), indicating that the height difference is attributable to other factors rather than mass being a confounding variable. 

Energy conservation during the flight period was tested to determine if height was a valid measure of energy output. As found in \cref{tab:ttest_paired}, the paired $t$-test within each group found no statistically significant difference between $KE$ at launch and $GPE$ at peak height (happy: $p = 0.41$, sad: $p = 0.43$). This confirms that mechanical energy was conserved during the flight period and that $GPE$ at peak reliably reflects $KE$ at launch. This also indicates that $EPE$ was successfully transferred into upward motion.

The force-compression measurement gave an $EPE$ of \qty{294.0}{\milli\joule}. As seen by \cref{tab:happy} and \cref{tab:sad}, mean $KE$ at launch for the happy poppers was \qty{0.0941}{\joule} and \qty{0.0814}{\joule} for the sad poppers, representing approximately \num{32}\% and \num{28}\% of the stored $EPE$, respectively. This suggests that approximately \num{70}\% of the elastic energy is lost during the initial compression phase, likely as heat, sound, or internal deformation. Once launched, however, both types conserve energy with comparable efficiency, further supporting the conclusion that facial expression does not influence energy transfer.

Overall, the statistical results and graphical evidence  support the hypothesis that facial expression does not influence maximum jump height. Instead, the motion of the poppers is governed primarily by elastic properties, mass, and energy conservation mechanisms rather than design features.

\subsection{Source of Experimental Error} 
A possible source of experimental error is the force at which each popper is launched. Despite the fixed location, the launch was performed by an individual rather than a repeatable machine; thus, the angle, speed, and pressure applied during compression may have varied, causing variability in the data. Additionally, error in human reaction and observation during video analysis may have occurred. Measurements taken from a meter stick in videos can be misread, leading to inconsistencies. Finally, material wear is a likely cause of error. As the poppers had been used previously, internal deformations in the rubber may have altered the elastic properties over time, leading to inconsistent energy results.

    
    
    
    
    
    
    
    
\section{Acknowledgements}
KC, SD, and SH primarily led data collection and recording. SH performed all popper compressions and authored the abstract. SY authored the introduction, methods and materials, generated the \texttt{R} code for graphical analysis, and revised the manuscript. RW calculated results based on height data. SD authored the sources of experimental error and performed all statistical analyses. KC performed the force-compression measurements, generated setup figures, and authored the discussion. All members contributed to proofreading, editing, and providing feedback on the final paper. We thank Dr. Evangelista for providing guidance, materials, and supervision throughout this laboratory exercise.








\bibliography{example.bib}
% References
% Tipler, P. A., & Mosca, G. (2003). Physics for Scientists and Engineers (5th ed.). W. H. Freeman.
% Chazot, C. (2025). FizziQ (Version 5.0.4) [Mobile application]. Trapèze. Available on Apple App Store.
% Google LLC. (2026). Google Sheets. Retrieved from https://www.google.com/sheets/about/
% R Core Team. (2025). R: A Language and Environment for Statistical Computing. R Foundation for Statistical Computing, Vienna, Austria. Retrieved from https://www.R-project.org/
% Wickham, H. (2016). ggplot2: Elegant Graphics for Data Analysis. Springer-Verlag New York. https://ggplot2.tidyverse.org


\end{document}
